Aging affects gait motion through well-documented biomechanical changes: reduced velocity, shorter stride length, compensatory hip strategies, and increased double-support time~\cite{boyer2023nia, studenski2011gait}. These kinematic signatures carry clinical significance and can predict mortality in older adults~\cite{studenski2011gait}.

Meanwhile, Spatial-Temporal Graph Convolutional Networks (ST-GCN)~\cite{yan2018stgcn} have shown great consistency in skeleton-based action recognition, and text-to-motion generative models such as MDM~\cite{tevet2023human} have been impressive in generating realistic motion from language descriptions. However, current models and datasets are biased heavily towards healthy young adults, and whether motion recognition and generation networks can grasp subtle changes caused by biomechanical factors remains unexplored.

We ask: \emph{Can age-specific kinematic signatures be extracted from clinical datasets, learned as a continuous representation, and reproduced by a generative model?}